\documentclass[12pt]{article}
\usepackage[utf8]{inputenc}
\usepackage[T1]{fontenc}
\usepackage{hyperref}
\title{Kamula Language Resources for K-12 Curriculum Development}
\author{Compiled by Claude Sonnet 5.5 (claude-sonnet-5-5) for Andr\'as Kornai}
\date{}
\begin{document}
\maketitle

\section{Introduction}
Kamula (ISO 639-3: xla; Glottocode: kamu1260), also called Kamira or Wawoi, is spoken in the northern Western Province of Papua New Guinea. It belongs to the small Kamula--Elevala family, together with Aekyom (Awin) and Pa (Pare). Sources give different speaker counts: about 800 (Usher, 1997), about 1,100 (Wikipedia, citing 2000) and about 2,000 (Joshua Project, citing Ethnologue 2016). Speakers live in two separate enclaves. The northern one is Keseki and Samokopa, near Wawoi Falls. The southern one is Wasapeya, on a lagoon north of the Aramia River. Very little published material exists, and I found no evidence of school materials in Kamula.

\section{Primary Educational Resources}
No primers, textbooks or school programmes for Kamula turned up in my searches. A mother-tongue curriculum would therefore have to be built from scratch, probably starting from the orthography and grammar work listed below.

\section{Linguistic and Reference Materials}
\begin{itemize}
\item Kamula language overview (Wikipedia). It cites Judy Routamaa's \emph{Kamula grammar essentials} (1994), \emph{Kamula phonology essentials} (1995) and an orthography paper (1997). I did not open these documents. \url{https://en.wikipedia.org/wiki/Kamula_language}
\item Kamula--Elevala languages overview. It cites Suter and Usher (2017), ``The Kamula-Elevala language family'', \emph{Language and Linguistics in Melanesia} 35. \url{https://en.wikipedia.org/wiki/Kamula\%E2\%80\%93Elevala_languages}
\item New Guinea World entry on Kamula (Timothy Usher). \url{https://newguineaworld.linguistik.uzh.ch/families/trans-new-guinea/central-west-new-guinea/digul-river-ok/kamula-elevala-river/kamula}
\item New Guinea World, Kamula--Elevala River family page. \url{https://newguineaworld.linguistik.uzh.ch/families/trans-new-guinea/central-west-new-guinea/digul-river-ok/kamula-elevala-river/start}
\end{itemize}
The search also mentioned, without a direct link, Iska and Judy Routamaa's \emph{Dialect survey report of the Kamula language, Western Province} (1996) and Reesink's \emph{Languages of the Aramia River Area} (1976). They are not listed as resources because I have no URL for them.

\section{Religious and Community Resources}
Joshua Project has a people-group page for the Kamura (Kamula). I did not verify its contents beyond the search result. \url{https://joshuaproject.net/people_groups/12481}

\section{Digital and Audio Resources}
I found no audio, video or corpus resources for Kamula.

\section{Curriculum Development Context}
Documentation is limited to survey and grammar sketches from the 1990s. The language is probably still spoken by all members of the community. A K-12 programme would need collaboration with the community and with the authors of the existing orthography and grammar work. Counts of speakers disagree between sources and should be checked locally.

\end{document}